\scanpage{92}\mspage{10}
% archive foliation: 91
\inserted{(où $t:S\longrightarrow U_{03}$)}
par la notation ordinaire
\[
(g,\lambda).\xi,
\]
et on trouve que l'on a, pour l'action
modulaire $T_{g,\lambda}$ sur $\xi$:
\[
(51)\qquad T_{g,\lambda}\xi
=\lambda\deleted{\unsure{\operatorname{sg}(g)}}(g,\lambda).\xi
\]
% margin position: right
\margin{$\xi\in\unsure{\Gamma(\underline{O}_{U_{0,3}}(1)_t)}$}
donc
\[
(52)\qquad T_{g,\lambda}\deleted{\unsure{\xi}}\inserted{\eta}^{\unsure{2|\deleted{2}}}
=\lambda^2(g,\lambda).\deleted{\unsure{\xi}}\inserted{\eta}^{\unsure{2|\deleted{2}}}
\]
% margin position: right
\margin{$\eta\in\unsure{\Gamma\underline{O}_{U_{0,3}}(2)(t)}$}
donc
\[
(53)\qquad T_{g,\lambda}(\tau)=\lambda^2\operatorname{sg}(g)\,(g,\lambda).\tau
\]
% margin position: right
\margin{$\tau\in\Gamma\deleted{\unsure{T_{U_{03}}}}\inserted{\unsure{\underline{O}_{U_{0,3}}(2)\otimes\omega_{S_0}(t)}}$}
Quand on se restreint au sous-groupe
$\mathbb{D}'_{3S_0}$, pour lequel on a $\lambda^2=\operatorname{sg}(g)$, on
trouve donc pour (51) (52) (53)
\[
(51\ \mathrm{bis})\qquad
T_{g\deleted{\unsure{,\lambda}}}\deleted{\unsure{\xi^2}}\inserted{\xi}
=\deleted{\unsure{\operatorname{sg}(g)}}\chi_4(g)\,(g.\xi)
\]
\[
(52\ \mathrm{bis})\qquad T_g\eta=\operatorname{sg}(g^0)\,(g.\eta)
\]
\[
(53\ \mathrm{bis})\qquad T_g(\tau)=\deleted{\unsure{\lambda}}g.\tau
\]
où $g\in\Gamma(\mathbb{D}'_{3S})$ pour $S$ sur $S_0$
\[
(54)\qquad \chi_4:\mathbb{D}'_{3S_0}\longrightarrow\unsure{\boldsymbol{\mu}|\mu}_4
\quad\text{le caractère canonique}
\]
\inserted{(de sorte que $\mathbb{D}'_{3\mathrm{ab}}\simeq\unsure{\boldsymbol{\mu}|\mu}_4$)}
